% TOP_PROBLEM: {"id": 9400078, "title": "Generalized Ramanujan conjecture for automorphic representations", "category_id": 1, "category": {"id": 1, "name": "number_theory", "display_name": "Number Theory", "description": "Properties of integers, prime numbers, Diophantine equations.", "slug": "number-theory", "order_index": 1, "created_at": "2026-07-31T15:26:25.670Z"}, "status": "partially_solved"}
% FIRST_POSTED: 2026-09-27
% !TeX program = lualatex
% ATTEMPT_STATUS: unresolved
% Research continuation by GPT_6_Astra_Ultra, 27 September 2026.
% Catalogue statements and existing sources preserved verbatim.
\documentclass[11pt]{article}
\usepackage[margin=25mm]{geometry}
\usepackage{fontspec}
\setmainfont{Latin Modern Roman}
\usepackage{amsmath,amssymb,mathtools}
\usepackage{unicode-math}
\setmathfont{Latin Modern Math}
\usepackage{xurl,hyperref}
\hypersetup{colorlinks=true,urlcolor=blue}
\setcounter{secnumdepth}{0}
\setlength{\parindent}{0pt}
\setlength{\parskip}{5pt}
\setlength{\emergencystretch}{3em}
\begin{document}
\section*{\texorpdfstring{Generalized Ramanujan conjecture for automorphic representations}{Generalized Ramanujan conjecture for automorphic representations}}\label{9400078}
\subsection{Definitions and mathematical statement}
Write a cuspidal automorphic representation as $\pi=\bigotimes'_v\pi_v$. A unitary irreducible local representation is tempered if it is weakly contained in the regular representation; for the groups in question this also has the usual matrix-coefficient characterization modulo the center. The conjecture is
\[
 \pi\text{ cuspidal on }\operatorname{GL}_n({\mathbb{A}}_F),\quad
 \omega_\pi\text{ unitary}\quad\Longrightarrow\quad
 \pi_v\text{ tempered for every place }v.
\]
At an unramified finite place, in unitary normalization, this asks that all Satake eigenvalues have absolute value one. The target includes ramified and archimedean places as well.
\par\smallskip\textit{Formulation and concept references:} \href{https://web.math.princeton.edu/~sarnak/Preprints/SarnakFieldsNotes.pdf}{[S1]}.

\subsection{Short English statement}
Must every local component of a unitary cuspidal automorphic representation satisfy the strongest expected temperedness bound?
\subsection{Research attempt}
\paragraph{Status, normalization, and attack.}
This is an unresolved attempt, dated 27 September 2026. It develops an exact
growth criterion at an unramified finite place, examines whether positivity
of Rankin--Selberg coefficients supplies it, and then checks why the criterion
cannot simply be copied to every place in the catalogue statement.
All Satake parameters below use the unitary normalization of [S1]. The
elementary power-sum argument is proved here; no new automorphic estimate
is claimed. The local representation facts used to interpret the criterion
are standard and are referenced explicitly.

Fix a cuspidal $\pi$ as in the statement and an unramified finite place $v$.
Write $q=\#(\mathcal O_F/v)$ and let
$\alpha_1,\ldots,\alpha_n\ne0$ be the Satake parameters. Unitarity of the
central character gives
\[
 \left|\prod_{j=1}^n\alpha_j\right|=1.
 \tag{30.1}
\]
Set $p_m=\sum_j\alpha_j^m$ and $R=\max_j|\alpha_j|$. The target here is
$R=1$: together with (30.1), the upper bound $R\leq1$ forces every
$|\alpha_j|=1$.

\paragraph{1. Recovering the largest parameter from traces.}
For an arbitrary finite nonzero complex multiset, one has
\[
 \limsup_{m\to\infty}|p_m|^{1/m}=R.
 \tag{30.2}
\]
The upper bound follows from $|p_m|\leq nR^m$. For the lower bound use the
generating function, initially for $|z|<1/R$,
\[
 F(z)=\sum_{m\geq1}p_mz^m
     =\sum_{j=1}^n\frac{\alpha_jz}{1-\alpha_jz}.
 \tag{30.3}
\]
Group together equal parameters. A distinct value $\beta$ of multiplicity
$r_\beta$ contributes $r_\beta\beta z/(1-\beta z)$, with a genuine pole at
$z=1/\beta$. All terms belonging to other distinct values are analytic
there, and $r_\beta$ is a positive integer. Thus no pole cancels. The nearest
pole has modulus $1/R$, which is the Taylor series radius of convergence.
The root test for power series proves (30.2).

This is stronger than examining a few traces: cancellations can be extensive
without changing the limiting growth. Equivalently, unramified temperedness
holds precisely when
\[
 \text{for every }\varepsilon>0\text{ there exists }C_{\varepsilon,v}
 \text{ such that }|p_m|\leq C_{\varepsilon,v}q^{\varepsilon m}
 \quad(m\geq1).
 \tag{30.4}
\]
The constants may depend on this representation and place. An estimate
with a single fixed positive exponent is insufficient. When temperedness
holds, the sharper bound $|p_m|\leq n$ is immediate; the point of (30.4)
is that it would suffice to prove only subexponential growth.

\paragraph{2. The positive local Rankin--Selberg criterion.}
Because $\pi_v$ is unitary, the Satake multiset of its contragredient agrees
with the complex conjugate multiset. Consequently the unramified factor is
\[
 L_v(s,\pi\times\widetilde\pi)
   =\prod_{i,j}(1-\alpha_i\overline{\alpha_j}q^{-s})^{-1},
\]
and in a right half-plane its logarithm has the absolutely convergent
expansion
\[
 \log L_v(s,\pi\times\widetilde\pi)
   =\sum_{m\geq1}\frac{|p_m|^2}{m}q^{-ms}.
 \tag{30.5}
\]
In particular every displayed coefficient is nonnegative. From (30.2),
\[
 \limsup_m\left(\frac{|p_m|^2}{m}\right)^{1/m}=R^2.
\]
It follows that temperedness is equivalent to the family of finite-energy
conditions
\[
 \mathcal E_v(\sigma):=
 \sum_{m\geq1}\frac{|p_m|^2}{m}q^{-m\sigma}<\infty
 \qquad\text{for every real }\sigma>0.
 \tag{30.6}
\]
Indeed convergence for every $\sigma>0$ implies
$R^2q^{-\sigma}\leq1$ for each $\sigma$, hence $R\leq1$; the converse
follows from $|p_m|\leq n$. Under (30.1), $R\geq1$, and the abscissa of
convergence of this local positive series is exactly
\[
 \sigma_v=2\log_q R.
 \tag{30.7}
\]
No assertion about convergence on the boundary is needed.

Positivity is useful because there is no cancellation between the terms in
(30.6). It also reveals the limitation of a global argument. Even if a
global positive majorant supplies convergence of this local series for
every $\sigma>1$, the root test only yields $R\leq q^{1/2}$. The desired
threshold is zero, not one. Sarnak's notes [S1] describe the stronger known
general bound $|\log_q|\alpha_j||\leq1/2-1/(n^2+1)$ and its origins.
That positive exponent is useful information but is not (30.4).
This reference is not an assertion of an exhaustive survey of later
improvements in special families.

Nor can meromorphic continuation of a global $L$-function be substituted for
convergence of its positive coefficient series. Analytic continuation changes
the domain on which a function is defined; it does not change the radius of
convergence of a particular series. For example $\sum_{m\geq1}R^{2m}z^m/m$
has its first singularity at $z=R^{-2}$ even though its expression
$-\log(1-R^2z)$ can be continued along many paths beyond that circle.
An argument using a continued value in place of the sum would lose the
positivity premise exactly where it is needed.

\paragraph{3. Stress tests for the proposed positivity attack.}
Consider the two parameters $(q^t,q^{-t})$, where $0<t<1/2$. They satisfy
the determinant condition and are the parameters of a spherical unitary
complementary-series representation of $\operatorname{GL}_2(F_v)$.
The existence of such local representations is part of the local unitary
picture discussed in [S1]. Here
\[
 p_m=q^{tm}+q^{-tm},\qquad
 |p_m|^2=q^{2tm}+2+q^{-2tm},
\]
so all the required positivity holds but $\sigma_v=2t>0$.
The corresponding local factor is
\[
 \frac{1}{(1-q^{2t-s})(1-q^{-s})^2(1-q^{-2t-s})}.
 \tag{30.8}
\]
Thus local unitarity, unitary determinant, and nonnegative logarithmic
coefficients do not prove temperedness. These are local test objects;
their occurrence in a cuspidal global representation is not being asserted.
Excluding that occurrence requires a genuinely global restriction.

A separate algebraic example tests cancellation: for $R>1$ take
$(R,-R,R^{-1},-R^{-1})$. Its product is one, all odd traces vanish, and
\[
 p_{2k}=2(R^{2k}+R^{-2k}).
 \tag{30.9}
\]
Half the traces therefore look as small as possible while (30.2) detects
the full growth. More generally multiplying both $R$ and $R^{-1}$ by every
$M$th root of unity makes all traces of orders $1,\ldots,M-1$ vanish.
The product of the resulting $2M$ parameters is one. This last family
changes the dimension with $M$ and is only a multiset example, not a
claimed family of automorphic representations. Its purpose is to show why
finite trace inspection cannot replace a quantified growth theorem.

\paragraph{4. Ramification requires the representation, not a copied formula.}
At a ramified finite place the standard local $L$-factor sees only the
inertia-invariant part of the monodromy kernel of a Weil--Deligne parameter.
It is not generally a degree-$n$ polynomial with all roots of modulus one
even when the representation is tempered. In the convention
$|\operatorname{Frob}_v|=q^{-1}$, consider the unitary Steinberg
representation of $\operatorname{GL}_2(F_v)$. Its Weil--Deligne data can be
written
\[
 r(\operatorname{Frob}_v)=
 \begin{pmatrix}q^{-1/2}&0\\0&q^{1/2}\end{pmatrix},\qquad
 Ne_2=e_1,\quad Ne_1=0,
\]
with trivial inertia. On $\ker N=\mathbb Ce_1$ the eigenvalue is
$q^{-1/2}$, so
\[
 L_v(s,\mathrm{St}_2)=(1-q^{-s-1/2})^{-1}.
 \tag{30.10}
\]
Nevertheless this representation is square integrable modulo its center,
hence tempered. The parameter normalization and the square-integrability
criterion are explained in [E1]. The matrix calculation verifies explicitly
why blindly imposing modulus one on the single standard-factor parameter
would reject a valid tempered representation.

For full local scope one may instead use the Langlands classification:
write the irreducible representation as the Langlands quotient of normalized
induction from tempered blocks $\tau_j|\det|^{t_j}$, with decreasing real
exponents. Temperedness requires all $t_j=0$ in this classification [E1].
That is a meaningful ramified target, but the proof of (30.2) for spherical
Satake traces has not bounded these block exponents. Additional twisted
local factors or matrix-coefficient estimates would need separate arguments.

At archimedean places there is no residue cardinality $q$ and no sequence
of Frobenius powers to which (30.6) applies. Even the spherical real
$\operatorname{GL}_2$ case contains complementary parameters $\pm t$;
in the usual weight-zero Laplace normalization they correspond to
$1/4-t^2$, whereas tempered principal-series parameters have zero real
part. The general statement also includes nonspherical archimedean
representations. A proof at all unramified finite places alone does not
prove the conjecture as stated.

\paragraph{Verification and outcome.}
The local companion script checks (30.9) exactly for $R=2$ through order
twelve. It also multiplies formal power series for (30.8) at $q=16,t=1/4$:
the logarithmic-derivative coefficients agree with
$(2^m+2^{-m})^2$ through order twelve, and its coefficients are nonnegative.
These checks verify normalization and bookkeeping, not automorphic
realizability. The Steinberg matrices satisfy
$r(\operatorname{Frob})Nr(\operatorname{Frob})^{-1}=q^{-1}N$, and restricting
to their monodromy kernel gives (30.10). The limiting statements are proved
by the rational generating function, not inferred from those finite tests.

The resulting conditional route is precise: establish (30.6) for every
unramified place, then separately force the Langlands block exponents to
vanish at all remaining places. No estimate doing the first task is proved
here, and the second task has not been reduced to the first. Positivity,
global continuation, and determinant normalization each survive explicit
tests that fail temperedness. The all-place generalized Ramanujan conjecture
therefore remains \textbf{UNRESOLVED}.

\subsection{Sources}
\begin{itemize}
\item[S1]Peter Sarnak, Notes on the Generalized Ramanujan Conjectures, in Harmonic Analysis, the Trace Formula, and Shimura Varieties.
\url{https://web.math.princeton.edu/~sarnak/Preprints/SarnakFieldsNotes.pdf}.
\textit{Formulation source.}
\item[S2]A note on Sarnak's density hypothesis for Sp4.
\url{https://ems.press/content/serial-article-files/53056}.
\textit{Further reference; not a proof endorsement.}
\item[S3]Regularized spectral expansion of a Rankin-Selberg period and moments.
\url{https://arxiv.org/abs/2609.01358}.
\textit{Further reference; not a proof endorsement.}

\item[E1]Torsten Wedhorn, \emph{The local Langlands correspondence for
$\operatorname{GL}(n)$ over $p$-adic fields}, arXiv:math/0011210v2
(25 November 2000), Sections 2.3, 3.2, and 4.3.
\url{https://arxiv.org/pdf/math/0011210v2}.
\textit{Author exposition of the local classification, temperedness,
Weil--Deligne factors, and the unitary Steinberg normalization. These are
established inputs, not new claims of this attempt.
Consulted 27 September 2026.}
\end{itemize}
\end{document}
